% ============================================================
% CAPITOLO 1 - INTRODUZIONE E CONTESTO
% ============================================================

\section{Introduzione}

\subsection{Il Problema della Comunicazione tra Sordi e Udenti}

Secondo l'Organizzazione Mondiale della Sanità (OMS) \cite{who2024deafness}, oltre il 5\% della popolazione mondiale --- circa 430 milioni di persone --- convive con una perdita uditiva disabilitante, di cui 34 milioni sono bambini. Le proiezioni indicano che entro il 2050 quasi 2,5 miliardi di individui presenteranno un qualche grado di perdita uditiva, e oltre 700 milioni necessiteranno di riabilitazione \cite{who2024deafness}. Il costo economico globale della perdita uditiva non trattata è stimato in circa 1 trilione di dollari all'anno, considerando spese sanitarie, supporto educativo, perdite di produttività e impatti sociali \cite{who2024hearingday}.

La \textbf{Lingua dei Segni} rappresenta il mezzo di comunicazione primario per la comunità sorda. Esistono oltre 300 lingue dei segni nel mondo, ciascuna con una propria grammatica, sintassi e struttura linguistica autonoma, distinta dalle corrispondenti lingue parlate \cite{wfd2024signlanguage}. La Lingua dei Segni Americana (ASL), utilizzata da circa mezzo milione di persone negli Stati Uniti e in Canada, è una delle più studiate e documentate.

Nonostante la sua ricchezza espressiva, la lingua dei segni crea una significativa \textbf{barriera comunicativa} tra la comunità sorda e la maggioranza della popolazione udente:
\begin{itemize}
    \item \textbf{Ambito sanitario:} i pazienti sordi affrontano frequentemente incomprensioni, ritardi diagnostici e difficoltà nell'accesso alle cure, a causa della carenza di interpreti professionisti e di personale sanitario competente nella lingua dei segni \cite{edf2023barriers}.
    \item \textbf{Ambito lavorativo:} il 62\% delle persone sorde e il 66\% delle persone con problemi d'udito riportano limitazioni nelle opportunità professionali dovute a barriere comunicative \cite{hamiltonrelay2023}.
    \item \textbf{Ambito educativo:} nei paesi in via di sviluppo, i bambini con perdita uditiva raramente ricevono un'istruzione adeguata, compromettendo le loro prospettive future \cite{who2024deafness}.
    \item \textbf{Vita quotidiana:} interazioni apparentemente semplici --- una telefonata, una conversazione informale, un'emergenza --- diventano ostacoli significativi in assenza di strumenti di mediazione.
\end{itemize}

In questo contesto, la tecnologia può svolgere un ruolo fondamentale nell'abbattere queste barriere, rendendo la comunicazione più accessibile e inclusiva.

% -----------------------------------------------------------

\subsection{Situazione Attuale nel Sign Language Recognition}

Il \textit{Sign Language Recognition} (SLR) è un campo di ricerca interdisciplinare che combina Computer Vision, Natural Language Processing e Deep Learning per tradurre automaticamente i gesti della lingua dei segni in testo o parlato.

\subsubsection{Approcci Tradizionali}

I primi sistemi di riconoscimento della lingua dei segni si basavano su \textbf{dispositivi indossabili}: guanti sensorizzati dotati di accelerometri, giroscopi e sensori di flessione, capaci di catturare con precisione la posizione e il movimento delle dita \cite{kumar2024slr}. Sebbene questi sistemi raggiungano accuratezze elevate (fino al 99,9\% per set limitati di 15 segni dinamici) \cite{kumar2024slr}, presentano limiti significativi:
\begin{itemize}
    \item \textbf{Invasività:} l'utente deve indossare hardware dedicato, riducendo la naturalezza della comunicazione.
    \item \textbf{Costo:} i dispositivi professionali hanno costi proibitivi per un'adozione su larga scala.
    \item \textbf{Scalabilità:} la calibrazione è specifica per ogni utente e ogni dispositivo.
\end{itemize}

\subsubsection{Approcci Vision-Based con Deep Learning}

L'evoluzione delle reti neurali profonde e la crescente potenza computazionale hanno permesso lo sviluppo di sistemi \textbf{vision-based}, che utilizzano semplici webcam o fotocamere per catturare i gesti, eliminando la necessità di hardware dedicato \cite{rastgoo2024slrsurvey}.

Le architetture principali impiegate nel campo sono:

\begin{itemize}
    \item \textbf{Convolutional Neural Networks (CNN):} utilizzate per l'estrazione di feature spaziali da singoli frame, efficaci nel riconoscimento di gesti statici come il \textit{fingerspelling} (alfabeto manuale) \cite{lecun2015deeplearning}. Varianti come MobileNetV2 \cite{sandler2018mobilenetv2} consentono l'inferenza su dispositivi a risorse limitate.

    \item \textbf{Recurrent Neural Networks (RNN/LSTM):} le reti \textit{Long Short-Term Memory} \cite{hochreiter1997lstm} eccellono nell'elaborazione di sequenze temporali, catturando la dinamica del movimento che caratterizza i segni. Le architetture ibride \textbf{CNN-LSTM} combinano l'estrazione spaziale con l'analisi temporale, raggiungendo accuratezze superiori al 99\% su dataset di segni isolati \cite{kumar2024slr}.

    \item \textbf{Transformer e Attention Mechanisms:} modelli più recenti basati su meccanismi di attenzione \cite{vaswani2017attention} stanno emergendo come alternativa promettente, in particolare per il riconoscimento continuo (\textit{Continuous Sign Language Recognition}, CSLR), superando le combinazioni CNN-RNN nel riconoscimento di sequenze gestuali complesse \cite{rastgoo2024slrsurvey}.

    \item \textbf{Object Detection (YOLO):} algoritmi come YOLOv8 \cite{ultralytics2023yolov8} vengono impiegati per la classificazione in tempo reale di lettere e gesti, offrendo un buon compromesso tra velocità e accuratezza. Sistemi basati su questa architettura hanno raggiunto il 98\% di accuratezza nel riconoscimento dell'alfabeto ASL \cite{alrashidi2024asl}.
\end{itemize}

\subsubsection{Strumenti e Framework Chiave}

L'ecosistema tecnologico attuale offre diversi strumenti fondamentali per lo sviluppo di sistemi SLR:

\begin{itemize}
    \item \textbf{MediaPipe Holistic (Google):} framework open-source per il tracciamento simultaneo di pose del corpo, mani e volto in tempo reale \cite{mediapipe2020}. Estrae fino a 543 landmark 3D per frame, fornendo una rappresentazione compatta e ricca del gesto. È diventato lo standard de facto per l'estrazione di feature nella ricerca SLR.

    \item \textbf{OpenCV:} libreria di Computer Vision \cite{opencv2000} utilizzata per l'acquisizione e il preprocessing del segnale video, la manipolazione dei frame e l'integrazione con modelli di inferenza.

    \item \textbf{TensorFlow/Keras:} framework di Deep Learning \cite{tensorflow2015} per la costruzione, l'addestramento e il deployment di modelli neurali. L'API sequenziale di Keras facilita la prototipazione rapida di architetture LSTM e CNN.

    \item \textbf{TensorFlow.js:} versione JavaScript di TensorFlow che consente l'esecuzione di modelli di Machine Learning direttamente nel browser \cite{tensorflowjs2018}, abilitando applicazioni web senza necessità di server dedicati.
\end{itemize}

% -----------------------------------------------------------

\subsection{Perché Signify}

Date le barriere comunicative descritte e il panorama tecnologico attuale, nasce l'esigenza di uno strumento che sia al contempo tecnicamente efficace e concretamente accessibile. \textbf{Signify --- Voices Through Hands} è la risposta a questa esigenza.

\subsubsection{La Visione del Progetto}

Signify nasce con un obiettivo chiaro: \textbf{dare voce ai gesti}. Il sistema traduce il segnale video catturato da una semplice webcam in testo (e potenzialmente in parlato) attraverso una pipeline che combina Computer Vision e Deep Learning:

\begin{center}
    \textit{Video Webcam} $\longrightarrow$ \textbf{MediaPipe} (Estrazione Landmark) $\longrightarrow$ \textbf{LSTM} (Classificazione Temporale) $\longrightarrow$ \textit{Output Testuale}
\end{center}

Il progetto si concentra sulla Lingua dei Segni Americana (ASL) come punto di partenza, utilizzando il dataset \textit{ASL Citizen} \cite{desai2024aslcitizen}, con l'ambizione di estendere il supporto ad altre lingue dei segni in futuro.

Dal punto di vista del Machine Learning, il problema viene modellato come un task di \textbf{Classificazione Multiclasse su Serie Temporali}:
\begin{itemize}
    \item \textbf{Input:} una sequenza di frame video $X = \{x_1, x_2, \dots, x_T\}$ contenenti un gesto.
    \item \textbf{Output:} una classe $y \in Y$ corrispondente alla parola (\textit{gloss}) eseguita.
    \item \textbf{Feature:} vettori di 258 dimensioni per frame (132 landmark della posa, 63 per ciascuna mano), estratti tramite MediaPipe Holistic \cite{mediapipe2020}.
\end{itemize}

\subsubsection{Accessibilità e Scelte Tecnologiche}

A differenza delle soluzioni commerciali chiuse --- come DeepASL \cite{fang2017deepasl}, SignLink o Sign-Speak --- che presentano licenze costose e API proprietarie, Signify adotta un approccio radicalmente diverso.

Il sistema è progettato con una filosofia \textit{``low-barrier''}: richiede unicamente una webcam standard e un computer con risorse computazionali modeste. Non sono necessari sensori di profondità, guanti sensorizzati o GPU dedicate per l'inferenza. Questa scelta progettuale si traduce in:
\begin{itemize}
    \item \textbf{Hardware minimo:} una webcam integrata è sufficiente per l'acquisizione video.
    \item \textbf{Deployment flessibile:} l'architettura è progettata per il funzionamento sia locale che via browser.
    \item \textbf{Stack tecnologico aperto:} MediaPipe \cite{mediapipe2020}, TensorFlow \cite{tensorflow2015}, OpenCV \cite{opencv2000} e Python sono tutti strumenti gratuiti e ben documentati.
\end{itemize}

\subsubsection{Un Progetto Open-Source}

Signify è rilasciato come progetto \textbf{completamente open-source}, con il codice sorgente disponibile pubblicamente su GitHub. Questa scelta non è casuale, ma risponde a un principio fondamentale: \textit{uno strumento nato per abbattere barriere non deve crearne di nuove}. L'approccio open-source garantisce che:
\begin{itemize}
    \item Chiunque possa utilizzare, modificare e distribuire il software gratuitamente.
    \item Le istituzioni educative possano integrarlo nei percorsi formativi senza costi di licenza.
    \item Le organizzazioni della comunità sorda possano adattarlo alle proprie esigenze specifiche.
    \item La comunità di sviluppatori possa contribuire al miglioramento continuo del sistema.
\end{itemize}

Nei capitoli successivi verranno descritti in dettaglio il dataset utilizzato, la pipeline di estrazione e preprocessamento dei dati, e l'architettura del modello di classificazione.